\section{Introduction and Background}

Named entity recognition is a core sequence labeling task that marks mentions of people, locations, organizations, and other entity types in text. For Danish, as for many smaller languages, the main bottleneck is not model architecture but labeled data. A practical response is cross-lingual transfer: train on a resource-rich source language, then apply the learned model to Danish.

The project proposal framed the central question: if one source language is closer to Danish than another, should that source transfer better? This report follows up with a full experimental comparison of English and Norwegian as source languages. English is the canonical high-resource transfer source; Norwegian is a North Germanic neighbor that shares substantial syntactic and lexical structure with Danish. Since both can be used with the same multilingual encoder (mBERT or XLM-R), the setup lets us test whether linguistic proximity translates into better downstream NER.

The experiments span five conditions: monolingual source-language baselines, zero-shot transfer to Danish, transfer followed by Danish fine-tuning, Danish-only models, and a Danish data-efficiency curve. We also compute typological distances between the languages using WALS syntactic feature vectors from \texttt{lang2vec} \citep{littell2017uriel}, add Swedish and German as additional comparison points, and run a paired bootstrap test on the core EN-vs-NO comparison to assess statistical significance.

The report is organized to mirror the proposal while giving a fuller account of the completed work: related work, research questions, data, linguistic distance analysis, methodology, results, analysis, discussion, and conclusion.
